% ===== 原文 pp.25--27（PDF pp.26--28）=====
% tr-03 片段 b：模型 $\mathrm{T}^{U(1)}_{N,d}$（(4.29)--(4.33)）与 $d\lessgtr N$ 的相结构

\para{模型 $\mathrm{T}^{U(1)}_{N,d}$}

考虑由两个正整数 $N$ 与 $d$ 标记的模型 $\mathrm{T}^{U(1)}_{N,d}$：
\begin{align*}
    G &= U(1),\\
    V &= \C(-d)\oplus\C(1)^{\oplus N}\ni(p,x_1,\ldots,x_N),\\
    W &= p\,f(x_1,\ldots,x_N),\\
    \tW &= -t\sigma.
\end{align*}
% [原文 p.26]
$f(x_1,\ldots,x_N)$ 是一个 $d$ 次多项式，它「一般」（generic）的含义是：若对所有 $i$ 都有 $\partial f/\partial x_i=0$，则 $x_1=\cdots=x_N=0$。R-荷在相差一个规范变换的意义下唯一：$R=(2-d\epsilon,\epsilon,\ldots,\epsilon)$，并且在 $J=\ee^{\pi\ii\epsilon}$ 下满足荷整性条件。

标量势为
\begin{equation}
    U = |-d\sigma p|^2 + |\sigma x|^2 + \frac{e^2}{2}\left(-d|p|^2+|x|^2-\zeta\right)^2
    + |f(x)|^2 + |p\,\dd f(x)|^2. \tag{4.29}
\end{equation}
% [Tutor Note: p.26 式 (4.29) 与 p.24 式 (4.23) 逐项对照即可核对：前两项是逐场的 $|\sigma\cdot(\text{场})|^2$（$p$ 的规范荷为 $-d$，故出现 $-d\sigma p$），第三项是 D-项 $\frac{e^2}{2}(-d|p|^2+|x|^2-\zeta)^2$（电荷加权、含 FI 参数），后两项来自 $W=pf$ 的 F-项 $|\partial W|^2$：$\partial_pW=f$、$\partial_{x_i}W=p\,\partial_if$。数值锚点：令 $p=x=0$，则 $U=\frac{e^2}{2}\zeta^2$，与 p.24 的「真空能量 $e^2\zeta^2/2$」一致；而 $\sigma=0$ 时 D-项与 F-项同时为零的条件正是随后讨论的 $|x|^2=\zeta$、$f(x)=0$、$p=0$。]
经典真空的集合以及规范对称性破缺的模式，取决于 FI 参数 $\zeta$ 为正、为负还是为零。

当 $\zeta>0$ 时，中间一项为零要求 $x$ 取非零值，从而完全破缺规范群。相应地，$\sigma$ 必须为零。又由于 $f(x)$ 是一般的，$|p\,\dd f(x)|^2=0$ 与 $x\neq0$ 一起要求 $p=0$。因此，经典真空的空间是满足 $|x|^2=\zeta$ 与 $f(x)=0$ 的 $x$ 的空间模去规范 $U(1)$ 作用。这恰好是 $\CP^{N-1}$ 中由 $f=0$ 定义的超曲面 $X_f$。所有横截于该真空流形的模式都是有质量的，于是理论在低能下约化为靶空间为 $X_f$ 的非线性 sigma 模型。Kähler 类与 B-场类由
\begin{equation}
    [\omega]\simeq\zeta H\in\mathrm{H}^2(X_f,\R),\qquad
    [B]\simeq[(\theta+\pi d)H]\in\mathrm{H}^2(X_f,\R/2\pi\Z). \tag{4.30}
\end{equation}
给出。$\theta\to\theta+\pi d$ 的平移来自积掉有质量的 $p$ 多重态 [12]。结果表明，这一关系可能还有进一步的量子修正，它们在 $\zeta\to+\infty$ 极限下是指数小的：$[\omega-\ii B]=[(t-d\pi\ii)H]+O(\ee^{-t})$。

当 $\zeta<0$ 时，中间一项为零迫使 $p$ 取非零值，从而把规范群破缺到由 $d$ 次单位根构成的子群 $\Z_d\subset U(1)$。相应地，$\sigma$ 必须为零。此外，由 $f(x)$ 的一般性，$|p\,\dd f(x)|^2=0$ 要求 $x$ 为零。于是真空流形是一个点：$p=\sqrt{-\zeta/d}=:\langle p\rangle$，$x=0$，$\sigma=0$，其稳定子为 $\Z_d$。这一次 $x$ 模式是无质量的（对 $d\ge3$），理论在低能下约化为带超势 $W=\langle p\rangle f(x)$ 的 $x$ 场模型，即在低能下为 Landau--Ginzburg 轨形 $(\C^N/\Z_d,W=\langle p\rangle f)$。

当 $\zeta=0$ 时，中间一项与第四项为零要求 $p=0$ 与 $x=0$。此时真空方程的其余部分对 $\sigma$ 不施加任何条件。也就是说，$\sigma$ 可以取任意值。真空流形是非紧空间 $\C$。这就是 Coulomb 分支。

总结起来，FI 参数空间 $\ii\fz^*\cong\R$ 有两个相——几何相 $\zeta>0$ 与 Landau--Ginzburg 相 $\zeta<0$——它们被相边界 $\zeta=0$ 隔开，在相边界处存在非紧的 Coulomb 分支。

FI 参数按 $\zeta'=\zeta+(N-d)\log(\mu'/\mu)$ 跑动，并且轴 $U(1)$ R-对称性是有异常的，除非 Calabi--Yau 条件 $d=N$ 得到满足。
% [原文 p.27]
在大 $\sigma$ 值处，有效扭超势为
\begin{equation}
    \tW_{\mathrm{eff}}(\sigma) = -t\sigma + d\sigma(\log(-d\sigma/\mu)-1) - N\sigma(\log(\sigma/\mu)-1), \tag{4.31}
\end{equation}
而 Coulomb 真空的方程是 $\partial_\sigma\tW_{\mathrm{eff}}\equiv 0 \bmod 2\pi\ii\Z$，或
\begin{equation}
    \sigma^{N-d} = (-d)^d\,\ee^{-t}\mu^{N-d}. \tag{4.32}
\end{equation}
% [Tutor Note: p.27 式 (4.32) 可由对 (4.31) 求导直接核对：$\partial_\sigma\tW_{\mathrm{eff}}=-t+d\log(-d\sigma/\mu)-N\log(\sigma/\mu)$，令其 $\equiv0 \bmod 2\pi\ii\Z$ 并取指数，得 $(-d\sigma/\mu)^d(\sigma/\mu)^{-N}\ee^{-t}=1$，即 $\sigma^{N-d}=(-d)^d\ee^{-t}\mu^{N-d}$。另外，$d=N$ 时 (4.32) 左边为 $\sigma^0=1$、右边为 $(-N)^N\ee^{-t}$，故「方程有解」只给出对 $t$ 的约束 $\ee^t=(-N)^N$——这正是紧随其后的判别轨迹 (4.33)。]
当 $d=N$ 时，我们得到一族超共形场论，其 $\hat c=N-2$，由 $t$ 以及 $f$ 的参数所参数化。由于方程 (4.32) 只有在 $\ee^t=(-N)^N$ 时才有解（此时 $\sigma\neq0$ 可任意），判别轨迹 $\Delta$ 是位于 $t\equiv N\log N+N\pi\ii$ 处的一个点，
\begin{equation}
    \mathfrak{M}_t = \C/2\pi\ii\Z - \{[N\log N + N\pi\ii]\}. \tag{4.33}
\end{equation}
我们看到，Calabi--Yau 流形 $X_f$ 上的非线性 sigma 模型与 Landau--Ginzburg 轨形 $(\C^N/\Z_N,f)$ 连续相连。特别地，这两个理论的拓扑 B-模型是等价的，而拓扑 A-模型由解析延拓相联系。这就是 \emph{CY/LG 对应}。在目前的模型中，(4.20) 里的包含关系都取等号：$\mathfrak{M}_t=\mathfrak{M}_A^0$，且 $\mathfrak{M}_W=\mathfrak{M}_B^0$。

当 $d<N$ 时，FI 参数 $\zeta$ 由正跑到负，使得 $\Lambda^{N-d}:=\ee^{-t}\mu^{N-d}$ 是 RG 不变量。在能标 $\mu$（满足 $e\gg\mu\gg|\Lambda|$）处，理论是 Fano 流形 $X_f$ 上的非线性 sigma 模型。在低于 $|\Lambda|$ 的能量下，sigma 模型的描述不再成立。存在 $(N-d)$ 个有质量间隙的 Coulomb 真空，位于 $\sigma^{N-d}=(-d)^d\Lambda^{N-d}$；以及在 $\sigma=0$ 处的一个 Higgs 分支理论（对 $d>1$），它是 Landau--Ginzburg 轨形 $(\C^N/\Z_d,f)$。当 $3\le d<N$ 时，该 Higgs 分支理论在红外极限下进一步流向一个 $\hat c=N(1-2/d)$ 的超共形场论。当 $d=2$ 时，对偶数（相应地，奇数）$N$，Higgs 分支理论有两个（相应地，一个）有质量间隙的超对称基态。

当 $d>N$ 时，FI 参数 $\zeta$ 由负跑到正，使得 $\Lambda^{d-N}:=\ee^{t}\mu^{d-N}$ 是 RG 不变量。该理论可看作对应于 Landau--Ginzburg 轨形 $(\C^N/\Z_d,f)$ 的、$\hat c=N(1-2/d)$ 的超共形场论被一个维数为 $2N/d$ 的相关算子扰动后的结果。存在 $(d-N)$ 个有质量间隙的 Coulomb 真空，位于 $\sigma^{d-N}=(-d)^{-d}\Lambda^{d-N}$；以及在 $\sigma=0$ 处的一个 Higgs 分支理论，它是一般型超曲面 $X_f$ 上的非线性 sigma 模型。
% [Tutor Note: p.27 的「$(N-d)$ 个 Coulomb 真空」与 (4.32) 自洽，可作为代数核对：$d<N$ 时 $\sigma^{N-d}=\text{const}$ 恰有 $N-d$ 个解；$d>N$ 时原文改写成 $\sigma^{d-N}=(-d)^{-d}\Lambda^{d-N}$，把它与 (4.32) 取倒数比较（$\sigma^{d-N}=(-d)^{-d}\ee^{t}\mu^{d-N}$，且 $\Lambda^{d-N}=\ee^{t}\mu^{d-N}$）恰好一致，故真空数 $d-N$ 也吻合。注意 $d=N$ 时 $\sigma$ 的方程退化为 $\sigma^0=1$，没有离散 Coulomb 真空，这与「$d=N$ 对应一族 $\hat c=N-2$ 的超共形场论」相符。]
